\documentclass[a4paper,11pt]{article}

\usepackage[margin=1.5cm]{geometry}
\usepackage{multicol}
\usepackage{listings}
\usepackage{xcolor}
\usepackage[T1]{fontenc}

\lstset{
  basicstyle=\ttfamily\small,
  keywordstyle=\color{blue},
  commentstyle=\color{gray},
  breaklines=true
}

\begin{document}

\begin{center}
    {\Large \textbf{Formulario CMake y CMakeLists}}\\
\end{center}

%%%%%%%%%%%%%%%%%%%%%%%%%%%%%%%%%%%%%%%%%%%%%%%%%%%%%%%%%%%%%%%

\section*{¿Qué es CMake?}

CMake es un \textbf{Build System Generator}.\\

NO compila directamente el código.\\

CMake:
\begin{itemize}
    \item Lee archivos \texttt{CMakeLists.txt}
    \item Genera archivos de build:
    \begin{itemize}
        \item Makefiles
        \item Ninja
        \item Visual Studio Projects
    \end{itemize}
    \item Configura compiladores, flags y librerías
\end{itemize}

%%%%%%%%%%%%%%%%%%%%%%%%%%%%%%%%%%%%%%%%%%%%%%%%%%%%%%%%%%%%%%%

\section*{Flujo General}

\begin{center}
\texttt{
CMake -> Build System -> Compiler -> Executable
}
\end{center}

Ejemplo:
\begin{itemize}
    \item CMake
    \item Ninja
    \item g++
    \item programa final
\end{itemize}

%%%%%%%%%%%%%%%%%%%%%%%%%%%%%%%%%%%%%%%%%%%%%%%%%%%%%%%%%%%%%%%

\section*{Build Stages}

\begin{center}
\begin{tabular}{| c | c |}
\hline
Stage & Descripción \\
\hline\hline
Configure & CMake analiza CMakeLists.txt \\
\hline
Generate & Genera Makefiles/Ninja/etc \\
\hline
Build & Compilación del código \\
\hline
Link & Une objetos y librerías \\
\hline
Run & Ejecuta el programa \\
\hline
\end{tabular}
\end{center}

%%%%%%%%%%%%%%%%%%%%%%%%%%%%%%%%%%%%%%%%%%%%%%%%%%%%%%%%%%%%%%%

\section*{Estructura básica}

\begin{lstlisting}
cmake_minimum_required(VERSION 3.20)

project(MyProject)

add_executable(app main.cpp)
\end{lstlisting}

%%%%%%%%%%%%%%%%%%%%%%%%%%%%%%%%%%%%%%%%%%%%%%%%%%%%%%%%%%%%%%%

\section*{Comandos básicos}

\begin{center}
\begin{tabular}{| c | c |}
\hline
Comando & Descripción \\
\hline\hline
project() & Nombre del proyecto \\
\hline
add\_executable() & Crear ejecutable \\
\hline
add\_library() & Crear librería \\
\hline
target\_link\_libraries() & Linkear librerías \\
\hline
target\_include\_directories() & Agregar includes \\
\hline
find\_package() & Buscar librerías instaladas \\
\hline
set() & Crear variables \\
\hline
message() & Imprimir mensajes \\
\hline
option() & Variables booleanas \\
\hline
install() & Instalación del proyecto \\
\hline
\end{tabular}
\end{center}

%%%%%%%%%%%%%%%%%%%%%%%%%%%%%%%%%%%%%%%%%%%%%%%%%%%%%%%%%%%%%%%

\section*{Compilación básica}

\begin{lstlisting}
cmake -B build
cmake --build build
\end{lstlisting}

Ejecutar:
\begin{lstlisting}
./build/app
\end{lstlisting}

%%%%%%%%%%%%%%%%%%%%%%%%%%%%%%%%%%%%%%%%%%%%%%%%%%%%%%%%%%%%%%%

\section*{Out-of-source builds}

Separar código fuente y archivos compilados.

\begin{lstlisting}
project/
|
|-- src/
|-- include/
|-- CMakeLists.txt
|-- build/
\end{lstlisting}

Ventajas:
\begin{itemize}
    \item Proyecto limpio
    \item Fácil borrar builds
    \item Soporte múltiples configuraciones
\end{itemize}

%%%%%%%%%%%%%%%%%%%%%%%%%%%%%%%%%%%%%%%%%%%%%%%%%%%%%%%%%%%%%%%

\section*{Versiones de C++}

\begin{lstlisting}
set(CMAKE_CXX_STANDARD 20)
set(CMAKE_CXX_STANDARD_REQUIRED ON)
\end{lstlisting}

%%%%%%%%%%%%%%%%%%%%%%%%%%%%%%%%%%%%%%%%%%%%%%%%%%%%%%%%%%%%%%%

\section*{Modo Debug y Release}

\begin{lstlisting}
set(CMAKE_BUILD_TYPE Debug)
\end{lstlisting}

Tipos:
\begin{itemize}
    \item Debug
    \item Release
    \item RelWithDebInfo
    \item MinSizeRel
\end{itemize}

%%%%%%%%%%%%%%%%%%%%%%%%%%%%%%%%%%%%%%%%%%%%%%%%%%%%%%%%%%%%%%%

\section*{Flags del compilador}

\begin{lstlisting}
target_compile_options(app PRIVATE -Wall -Wextra)
\end{lstlisting}

Ejemplos:
\begin{itemize}
    \item -Wall
    \item -Wextra
    \item -O2
    \item -g
\end{itemize}

%%%%%%%%%%%%%%%%%%%%%%%%%%%%%%%%%%%%%%%%%%%%%%%%%%%%%%%%%%%%%%%

\section*{Agregar múltiples archivos}

\begin{lstlisting}
add_executable(app
    src/main.cpp
    src/math.cpp
    src/utils.cpp
)
\end{lstlisting}

%%%%%%%%%%%%%%%%%%%%%%%%%%%%%%%%%%%%%%%%%%%%%%%%%%%%%%%%%%%%%%%

\section*{GLOB automático}

\begin{lstlisting}
file(GLOB_RECURSE SRC_FILES
    src/*.cpp
)

add_executable(app ${SRC_FILES})
\end{lstlisting}

%%%%%%%%%%%%%%%%%%%%%%%%%%%%%%%%%%%%%%%%%%%%%%%%%%%%%%%%%%%%%%%

\section*{Include directories}

\begin{lstlisting}
target_include_directories(app PRIVATE include/)
\end{lstlisting}

Tipos:
\begin{itemize}
    \item PRIVATE
    \item PUBLIC
    \item INTERFACE
\end{itemize}

%%%%%%%%%%%%%%%%%%%%%%%%%%%%%%%%%%%%%%%%%%%%%%%%%%%%%%%%%%%%%%%

\section*{Librerías}

\subsection*{Static Library}

\begin{lstlisting}
add_library(math STATIC math.cpp)
\end{lstlisting}

\subsection*{Shared Library}

\begin{lstlisting}
add_library(math SHARED math.cpp)
\end{lstlisting}

%%%%%%%%%%%%%%%%%%%%%%%%%%%%%%%%%%%%%%%%%%%%%%%%%%%%%%%%%%%%%%%

\section*{Linkear librerías}

\begin{lstlisting}
target_link_libraries(app PRIVATE math)
\end{lstlisting}

%%%%%%%%%%%%%%%%%%%%%%%%%%%%%%%%%%%%%%%%%%%%%%%%%%%%%%%%%%%%%%%

\section*{find\_package}

Buscar librerías instaladas:

\begin{lstlisting}
find_package(OpenCV REQUIRED)

target_link_libraries(app
    PRIVATE
    OpenCV::OpenCV
)
\end{lstlisting}

%%%%%%%%%%%%%%%%%%%%%%%%%%%%%%%%%%%%%%%%%%%%%%%%%%%%%%%%%%%%%%%

\section*{Toolchains}

Un Toolchain es el conjunto de herramientas usado para compilar:

\begin{itemize}
    \item Compiler
    \item Linker
    \item Assembler
    \item Sysroot
    \item Librerías
\end{itemize}

Ejemplos:
\begin{itemize}
    \item GCC
    \item Clang
    \item MSVC
    \item ARM Embedded Toolchain
\end{itemize}

%%%%%%%%%%%%%%%%%%%%%%%%%%%%%%%%%%%%%%%%%%%%%%%%%%%%%%%%%%%%%%%

\section*{Toolchain File}

Archivo usado para Cross Compilation.

\begin{lstlisting}
cmake -B build \
-DCMAKE_TOOLCHAIN_FILE=toolchain.cmake
\end{lstlisting}

Ejemplo:

\begin{lstlisting}
set(CMAKE_SYSTEM_NAME Linux)

set(CMAKE_C_COMPILER arm-linux-gcc)
set(CMAKE_CXX_COMPILER arm-linux-g++)
\end{lstlisting}

%%%%%%%%%%%%%%%%%%%%%%%%%%%%%%%%%%%%%%%%%%%%%%%%%%%%%%%%%%%%%%%

\section*{Cross Compilation}

Compilar en una arquitectura para otra.

Ejemplo:
\begin{itemize}
    \item PC x86
    \item Raspberry Pi ARM
\end{itemize}

%%%%%%%%%%%%%%%%%%%%%%%%%%%%%%%%%%%%%%%%%%%%%%%%%%%%%%%%%%%%%%%

\section*{Generadores}

\begin{center}
\begin{tabular}{| c | c |}
\hline
Generator & Plataforma \\
\hline\hline
Unix Makefiles & Linux \\
\hline
Ninja & Multiplataforma \\
\hline
Visual Studio & Windows \\
\hline
Xcode & macOS \\
\hline
\end{tabular}
\end{center}

Seleccionar generator:

\begin{lstlisting}
cmake -G Ninja -B build
\end{lstlisting}

%%%%%%%%%%%%%%%%%%%%%%%%%%%%%%%%%%%%%%%%%%%%%%%%%%%%%%%%%%%%%%%

\section*{Ninja}

Sistema de build rápido y moderno.

\begin{lstlisting}
sudo apt install ninja-build
\end{lstlisting}

%%%%%%%%%%%%%%%%%%%%%%%%%%%%%%%%%%%%%%%%%%%%%%%%%%%%%%%%%%%%%%%

\section*{Variables importantes}

\begin{center}
\begin{tabular}{| c | c |}
\hline
Variable & Descripción \\
\hline\hline
CMAKE\_SOURCE\_DIR & Directorio raíz \\
\hline
CMAKE\_BINARY\_DIR & Directorio build \\
\hline
CMAKE\_CXX\_COMPILER & Compilador C++ \\
\hline
CMAKE\_BUILD\_TYPE & Debug/Release \\
\hline
\end{tabular}
\end{center}

%%%%%%%%%%%%%%%%%%%%%%%%%%%%%%%%%%%%%%%%%%%%%%%%%%%%%%%%%%%%%%%

\section*{Mensajes}

\begin{lstlisting}
message("Hola")
message(STATUS "Configurando")
message(WARNING "Advertencia")
message(FATAL_ERROR "Error")
\end{lstlisting}

%%%%%%%%%%%%%%%%%%%%%%%%%%%%%%%%%%%%%%%%%%%%%%%%%%%%%%%%%%%%%%%

\section*{Testing}

\begin{lstlisting}
enable_testing()

add_test(NAME test1 COMMAND app)
\end{lstlisting}

Ejecutar tests:
\begin{lstlisting}
ctest
\end{lstlisting}

%%%%%%%%%%%%%%%%%%%%%%%%%%%%%%%%%%%%%%%%%%%%%%%%%%%%%%%%%%%%%%%

\section*{Instalación}

\begin{lstlisting}
install(TARGETS app
    DESTINATION bin
)
\end{lstlisting}

%%%%%%%%%%%%%%%%%%%%%%%%%%%%%%%%%%%%%%%%%%%%%%%%%%%%%%%%%%%%%%%

\section*{CMake Presets}

Archivo:
\begin{lstlisting}
CMakePresets.json
\end{lstlisting}

Ejemplo:
\begin{lstlisting}
{
  "version": 3,
  "configurePresets": [
    {
      "name": "debug",
      "binaryDir": "build",
      "cacheVariables": {
        "CMAKE_BUILD_TYPE": "Debug"
      }
    }
  ]
}
\end{lstlisting}

Usar preset:
\begin{lstlisting}
cmake --preset debug
\end{lstlisting}

%%%%%%%%%%%%%%%%%%%%%%%%%%%%%%%%%%%%%%%%%%%%%%%%%%%%%%%%%%%%%%%

\section*{Debug de CMake}

\begin{lstlisting}
cmake --trace
cmake --debug-output
\end{lstlisting}

%%%%%%%%%%%%%%%%%%%%%%%%%%%%%%%%%%%%%%%%%%%%%%%%%%%%%%%%%%%%%%%

\section*{Estructura moderna recomendada}

\begin{lstlisting}
project/
|
|-- src/
|-- include/
|-- tests/
|-- external/
|-- CMakeLists.txt
|-- CMakePresets.json
\end{lstlisting}

%%%%%%%%%%%%%%%%%%%%%%%%%%%%%%%%%%%%%%%%%%%%%%%%%%%%%%%%%%%%%%%

\section*{Flujo moderno recomendado}

\begin{lstlisting}
cmake -G Ninja -B build

cmake --build build

./build/app
\end{lstlisting}

%%%%%%%%%%%%%%%%%%%%%%%%%%%%%%%%%%%%%%%%%%%%%%%%%%%%%%%%%%%%%%%

\section*{Integración con IDEs}

\begin{itemize}
    \item VSCode
    \item CLion
    \item Visual Studio
    \item Qt Creator
\end{itemize}

Los IDEs normalmente:
\begin{itemize}
    \item Ejecutan CMake automáticamente
    \item Configuran Debuggers
    \item Detectan cambios
\end{itemize}

%%%%%%%%%%%%%%%%%%%%%%%%%%%%%%%%%%%%%%%%%%%%%%%%%%%%%%%%%%%%%%%


\section*{Ejemplo de CMakeLists Complejo}

Estructura del proyecto:

\begin{lstlisting}
project/
|
|-- src/
|   |-- main.cpp
|   |-- math/
|       |-- math.cpp
|
|-- include/
|   |-- math/
|       |-- math.hpp
|
|-- tests/
|
|-- external/
|
|-- CMakeLists.txt
\end{lstlisting}

%%%%%%%%%%%%%%%%%%%%%%%%%%%%%%%%%%%%%%%%%%%%%%%%%%%%%%%%%%%%%%%

\begin{lstlisting}
##############################################################
# Version minima requerida de CMake
#
# IMPORTANTE:
# Algunas features modernas solo existen
# en versiones recientes.
##############################################################

cmake_minimum_required(VERSION 3.20)

##############################################################
# Nombre del proyecto
#
# Tambien define variables internas:
#
# PROJECT_NAME
# PROJECT_SOURCE_DIR
#
# Puede incluir:
# VERSION
# DESCRIPTION
# LANGUAGES
##############################################################

project(
    ComplexProject
    VERSION 1.0
    DESCRIPTION "Ejemplo avanzado de CMake"
    LANGUAGES CXX
)

##############################################################
# Definir el standard de C++
#
# REQUIRED ON:
# Obliga a usar exactamente ese standard.
#
# Si OFF:
# Puede degradarse a versiones anteriores.
##############################################################

set(CMAKE_CXX_STANDARD 20)
set(CMAKE_CXX_STANDARD_REQUIRED ON)

##############################################################
# Export compile commands
#
# Genera compile_commands.json
#
# Util para:
# - clangd
# - LSPs
# - herramientas de analisis
##############################################################

set(CMAKE_EXPORT_COMPILE_COMMANDS ON)

##############################################################
# Tipo de build
#
# Debug:
# - simbolos de debug
# - sin optimizacion
#
# Release:
# - optimizaciones
#
# RelWithDebInfo:
# - optimizado + debug
##############################################################

if(NOT CMAKE_BUILD_TYPE)
    set(CMAKE_BUILD_TYPE Debug)
endif()

##############################################################
# Buscar todos los archivos cpp
#
# GLOB_RECURSE:
# busca recursivamente.
#
# Ventaja:
# no agregar archivos manualmente.
#
# Desventaja:
# algunos proyectos prefieren declarar
# explicitamente todos los archivos.
##############################################################

file(
    GLOB_RECURSE
    SRC_FILES
    src/*.cpp
)

##############################################################
# Crear libreria propia
#
# STATIC:
# libreria estatica.
#
# Otras opciones:
#
# SHARED:
# libreria dinamica (.so/.dll)
#
# INTERFACE:
# libreria solo headers.
##############################################################

add_library(
    mathlib
    STATIC
    src/math/math.cpp
)

##############################################################
# Includes de la libreria
#
# PUBLIC:
# targets que usen mathlib
# heredaran este include.
#
# PRIVATE:
# solo mathlib lo usa.
#
# INTERFACE:
# solo consumidores lo usan.
##############################################################

target_include_directories(
    mathlib
    PUBLIC
    include/
)

##############################################################
# Crear ejecutable principal
##############################################################

add_executable(
    app
    src/main.cpp
)

##############################################################
# Linkear libreria propia
#
# app depende de mathlib.
#
# CMake automaticamente:
# - ordena compilacion
# - agrega linker flags
# - propaga includes PUBLIC
##############################################################

target_link_libraries(
    app
    PRIVATE
    mathlib
)

##############################################################
# Warnings del compilador
#
# GNU/Clang:
# -Wall
# -Wextra
#
# MSVC usa otros flags.
##############################################################

target_compile_options(
    app
    PRIVATE
    -Wall
    -Wextra
)

##############################################################
# Buscar libreria externa
#
# REQUIRED:
# si no existe -> error.
#
# Otras librerias:
# - OpenCV
# - Boost
# - Eigen3
# - fmt
##############################################################

find_package(fmt REQUIRED)

##############################################################
# Linkear libreria externa
#
# fmt::fmt es un target exportado
# por la libreria.
#
# IMPORTANTE:
# modern CMake prefiere targets
# en lugar de paths manuales.
##############################################################

target_link_libraries(
    app
    PRIVATE
    fmt::fmt
)

##############################################################
# Definir macros de preprocesador
#
# DEBUG_MODE quedara disponible:
#
# #ifdef DEBUG_MODE
##############################################################

target_compile_definitions(
    app
    PRIVATE
    DEBUG_MODE
)

##############################################################
# Testing
#
# Habilita soporte para ctest.
##############################################################

enable_testing()

##############################################################
# Crear ejecutable de tests
##############################################################

add_executable(
    tests
    tests/test_main.cpp
)

##############################################################
# Linkear tests con libreria
##############################################################

target_link_libraries(
    tests
    PRIVATE
    mathlib
)

##############################################################
# Registrar test
#
# ctest ejecutara este binario.
##############################################################

add_test(
    NAME math_test
    COMMAND tests
)

##############################################################
# Instalacion
#
# DESTINATION:
# donde se instalara.
#
# Ejemplo:
# /usr/local/bin
##############################################################

install(
    TARGETS app
    DESTINATION bin
)

##############################################################
# Mensajes informativos
##############################################################

message(STATUS "Build type: ${CMAKE_BUILD_TYPE}")
message(STATUS "Compiler: ${CMAKE_CXX_COMPILER}")

##############################################################
# Diferentes compiladores
#
# GNU:
# gcc/g++
#
# Clang:
# clang/clang++
#
# MSVC:
# Visual Studio
##############################################################

if(CMAKE_CXX_COMPILER_ID STREQUAL "GNU")

    message(STATUS "Usando GCC")

elseif(CMAKE_CXX_COMPILER_ID STREQUAL "Clang")

    message(STATUS "Usando Clang")

elseif(CMAKE_CXX_COMPILER_ID STREQUAL "MSVC")

    message(STATUS "Usando MSVC")

endif()

##############################################################
# Cross compilation
#
# Si se usa:
#
# cmake -DCMAKE_TOOLCHAIN_FILE=...
#
# entonces CMake puede compilar
# para otra arquitectura.
#
# Ejemplo:
# x86 -> ARM
##############################################################

message(
    STATUS
    "System processor: ${CMAKE_SYSTEM_PROCESSOR}"
)

##############################################################
# Custom commands
#
# Ejecutar comandos personalizados.
#
# Ejemplo:
# copiar assets
# generar codigo
# ejecutar scripts
##############################################################

add_custom_command(
    TARGET app
    POST_BUILD

    COMMAND ${CMAKE_COMMAND} -E echo
    "Build finalizado"
)

##############################################################
# Custom targets
#
# Targets virtuales.
#
# Ejemplo:
# make clean_extra
##############################################################

add_custom_target(

    print_info ALL

    COMMAND ${CMAKE_COMMAND} -E echo
    "Proyecto compilado correctamente"

)

##############################################################
# Subdirectorios
#
# Util para proyectos grandes.
#
# Cada carpeta puede tener
# su propio CMakeLists.txt
##############################################################

# add_subdirectory(external)
# add_subdirectory(src/math)

##############################################################
# Presets modernos
#
# Se configuran en:
#
# CMakePresets.json
#
# Facilitan:
# - debug
# - release
# - cross compile
# - CI/CD
##############################################################
\end{lstlisting}

\end{document}